\documentclass[11pt]{article}

\usepackage[margin=1.2in]{geometry}
\usepackage{amsmath,amssymb}
\usepackage{microtype}
\usepackage{booktabs}
\usepackage{xcolor}
\usepackage{mdframed}
\usepackage[colorlinks=true,linkcolor=blue!50!black,urlcolor=blue!50!black]{hyperref}

\newmdenv[backgroundcolor=blue!4,linecolor=blue!40,linewidth=1pt,
  skipabove=8pt,skipbelow=8pt]{keyidea}

\title{The 3-adic Shadow of a Collatz Orbit:\\
       \Large A Friendly Guide}
\author{Companion notes to ``Divergence as Sustained
Compression''\thanks{Written in collaboration with Claude (Anthropic).
These notes explain the ideas; the paper and the Lean files contain the
proofs. The Collatz conjecture remains open.}}
\date{June 12, 2026}

\begin{document}
\maketitle

\section{Where we left off}

The first guide in this series ended with a picture. Playing the
Collatz game (halve if even; triple-and-add-one if odd, then halve),
a number can escape to infinity only by keeping the share of its
``odd'' steps above the critical $63.09\%$ forever --- and we know,
provably, that almost every number fails to do this almost
immediately, while no argument that examines numbers through any
bounded lens can ever decide who succeeds. The question this paper
asks is simple to state:

\begin{center}
\emph{Suppose a number did keep it up forever. What else would have to
be true about it?}
\end{center}

The answer turns out to be startlingly concrete, and it lives in an
unexpected place: not in the number's binary digits (base 2), which
the first paper exhausted, but in its \textbf{base-3 digits}.

\section{A magic cancellation}

Run a number $n$ through $T$ steps of the game, of which $j$ were odd
steps, and call the result $n_T$. There is an exact bookkeeping
identity --- every multiplication and every carry accounted for, no
approximations:
\[
2^T \cdot n_T \;=\; 3^{j} \cdot n + d,
\]
where the leftover term $d$ is built up step by step along the way.
Two facts about $d$ make everything happen.

First: $d$ \textbf{does not depend on $n$}. It depends only on the
\emph{pattern} of odd and even steps --- the word of choices the orbit
made. Two starting numbers with the same first-$T$ pattern have
exactly the same $d$. (Patterns and starting numbers are in
one-to-one correspondence with remainders mod $2^T$ --- that was the
parity bijection of the first paper.)

Second, the magic: look at the identity \emph{modulo} $3^j$, i.e.\
ask only for remainders upon division by $3^j$. The term
$3^j \cdot n$ is divisible by $3^j$ --- it vanishes --- taking the
only mention of $n$ with it:
\[
n_T \;\equiv\; (\text{something built from the pattern alone})
\pmod{3^{j}}.
\]

\begin{keyidea}
\textbf{The shadow.} After $T$ steps, the orbit's remainder modulo
$3^j$ --- its ``base-3 shadow'' --- is completely determined by the
odd/even pattern. The starting number has no say. (Verified
mechanically across twenty thousand random starts: no exceptions, as
the theorem demands.)
\end{keyidea}

\section{Minting coins with a small ledger}

Now count, because the two sides of this equation are counted in
different currencies.

\emph{How many shadows could there be?} The shadow lives among the
remainders modulo $3^j$, and there are $3^j$ of those. For a
surviving orbit, $j$ is locked to about $63.09\%$ of $T$ by the
critical-line results --- and here sits the coincidence that powers
the whole paper. Measured in bits, $3^j$ amounts to
$j \cdot \log_2 3 \approx 0.6309 \times 1.585 \times T$ bits, and
\[
0.6309\ldots \times 1.585\ldots = 1 \quad \text{exactly.}
\]
A surviving orbit mints \textbf{exactly one bit of base-3 information
per step}. Not approximately --- the two famous constants of this
problem multiply to precisely $1$.

\emph{How many patterns could be doing the minting?} Patterns with
$63\%$ odd steps are rare: coin-flip counting (Pascal's triangle
again) says there are only about $2^{0.95\,T}$ of them ---
$0.95$ bits per step of pattern.

So the mint produces $1.00$ bits per step of base-3 currency, printed
from a ledger containing only $0.95$ bits per step of instructions.
The currency must be \emph{redundant}: out of the $3^j$ possible
shadows, surviving orbits can only ever land on an exponentially tiny
fraction --- about $2^{-0.045\,T}$ of them. At a window of just $18$
steps this is already visible by brute force: surviving orbits hit
only $2.17\%$ of the available remainders modulo $3^{12}$.

\begin{keyidea}
\textbf{The reduction.} A number that escapes to infinity must emit
base-3 digits that are about $4.5\%$ compressible --- carrying
provably less information than their length --- at every scale, by
consistent amounts, forever. ``Divergence'' and ``a number that
out-compresses its own arithmetic'' are the same thing.
\end{keyidea}

Both halves of this are now machine-verified theorems: the density
floor (escapees are pinned to the critical $63.09\%$, where the mint
runs at exactly $1$ bit per step) and the confinement count (the
shadow set is small, with --- pleasingly --- the \emph{same} constants
$17/27$ and $3^T$ that ran the ``almost all numbers fall'' theorem in
the first paper; the base-2 and base-3 mirrors are one phenomenon).

\section{The ledger has a grammar: scales talk to each other}

The leftover term $d$ obeys a beautiful composition rule. If you run
a window of length $S$ and then a window of length $T$, with patterns
$w_1$ and $w_2$,
\[
d(w_1 w_2) \;=\; 3^{j_2}\, d(w_1) \;+\; 2^{S}\, d(w_2):
\]
the second window's base-3 weight scales the first ledger entry, the
first window's base-2 length scales the second. (Readers who know some
dynamics will recognize a skew product of the doubling and tripling
worlds; everyone else can read it as: \emph{the books at different
scales are not independent --- they're nested}.)

The experiments confirm this with force. At a window of $18$ steps,
knowing an orbit's full-window shadow pins down its half-window shadow
almost uniquely (about $1.09$ candidates on average, out of hundreds
possible). A diverging number wouldn't just need a compressible shadow
at each scale separately --- it would need all its shadows, at all
scales, to \emph{cohere} through this grammar. That consistency
requirement is a constraint nobody has ever exploited, and it is the
paper's first recommendation for future attack.

One more experimental finding sharpens the hunt. The slowest possible
escapees --- those hugging the critical line --- are forced into
extremely regular patterns (so-called Sturmian words, the most orderly
aperiodic sequences there are). Feeding such patterns through the
shadow map, they cluster dramatically: $400$ of them produced only
$167$ distinct shadows, where $400$ random patterns of the same
density produced $400$. The most dangerous escape routes live in the
\emph{thinnest} part of the base-3 world --- precisely where
arithmetic has the most leverage to forbid them.

\section{So is it solved?}

No --- and the paper is explicit about what stands in the way. The
remaining statement has a familiar shape: it asks, in essence, whether
a whole number's base-2 behavior and base-3 behavior can conspire ---
whether an integer can have such special joint structure in the two
bases that it sustains the compression. Mathematics has a famous
circle of problems about exactly this (``$\times 2$, $\times 3$''
problems, after Furstenberg), and they are notoriously hard: even the
innocent-sounding claim \emph{``the binary digits of $3^k$ look
random''} is open. Two of our own earlier results explain why no
shortcut exists: every finite behavior really occurs (so the
compressed shadow sets are never empty --- counting alone cannot
finish), and no bounded-memory argument can decide who inhabits them.

What the paper contributes is the conversion. Before: ``prove no
number escapes,'' a problem with no agreed-upon shape. After: a
single quantitative statement --- \emph{no integer compresses its own
base-3 digit stream by $4.5\%$ at every scale consistently forever}
--- with machine-checked constants, a numeric gap between $0.95$ and
$1.00$ that does not depend on any estimate being sharp, three
identified levers (cross-scale consistency, structured-pattern
exclusion, and the ergodic question for the rigidity specialists), and
an honest map of which lever sits in known-hard territory.

\section{Summary in four sentences}

An exact identity makes an orbit's remainder modulo $3^j$ a function
of its odd/even pattern alone --- the starting number cancels out.
Escaping numbers are pinned to the critical density, where they mint
exactly one bit of base-3 information per step from patterns worth
only $0.95$ bits per step --- so their base-3 digits must be
compressible by about $4.5\%$, at every scale, forever, and the
shadows at different scales are locked together by an exact
composition law. All of this is machine-verified; the experiments show
the true compression is even stronger, the scales nearly determine
each other, and the most regular escape patterns cluster in the
thinnest shadow sets. The Collatz conjecture remains open --- but its
divergence half now has a precise shape: it asks whether an integer
can beat the entropy of its own arithmetic.

\end{document}
